\section{Convergence of the free precision horizon}
\label{sec:free-horizon-convergence}
We prove the retained-horizon convergence used for the stationary paths.
The cutoff tends to infinity while the scale and block side remain fixed.
All coefficients are kept on the lattice cover. Finite-torus matrices are
used to identify their limits; they are not substituted for the cover
coefficients.

\subsection{Compactness with an exponential row bound}
A row $K$ between finite coordinate sets $I,J$ has coefficients
$K_{ab}(z)$, $z\in\mathbb Z^4$. Suppose, eventually in $N$,
\begin{equation}\label{eq:free-horizon-row-budget}
 \sum_z e^{\sigma|z|_1}|K_{N,ab}(z)|\le C
 \qquad(a\in I,b\in J),\qquad \sigma>0.
\end{equation}
A subsequence of any given cutoff sequence converges in every coefficient:
enumerate the countably many entries, extract successively in their bounded
real intervals, and take the diagonal subsequence. If its limit is $K$,
pass to the limit in each finite subsum of
\eqref{eq:free-horizon-row-budget}. The nonnegative sums remain bounded
by $C$, so the full sum has the same bound.

For every $0\le\tau<\sigma$ this coefficient convergence gives
\begin{equation}\label{eq:free-horizon-narrower}
       \sum_z e^{\tau|z|_1}|K_{N,ab}(z)-K_{ab}(z)|\longrightarrow0.
\end{equation}
Indeed its summands are at most
$2C e^{-(\sigma-\tau)|z|_1}$, and
\[
 \sum_{z\in\mathbb Z^4}e^{-t|z|_1}
       =\left(\frac{1+e^{-t}}{1-e^{-t}}\right)^4<\infty
       \qquad(t>0).
\]
First make the tail of this majorant small, and then use convergence
on its finite complement. This also proves convergence of the Fourier
symbols on real momenta and of the periodized matrices on every fixed
torus, since their entries are sums of subfamilies of the same coefficients.
The argument applies to a cross row by including its finitely many block
residues in $I$.

Two ways of identifying the resulting cluster limits will be useful.
If the Fourier symbols have a prescribed limit on a dense subset of
the real momentum torus, any two cluster rows have the same symbol
there. Their symbols are continuous by absolute summability, so they
agree everywhere. Integrating against $e^{-ip\cdot z}$ recovers the
coefficient $K_{ab}(z)$, proving equality of the rows. Alternatively,
it suffices that the cluster rows have equal matrices on an unbounded
sequence of torus sides. For a fixed displacement $z$ the appropriate
periodized entry is
\[
                       \sum_{h\in\mathbb Z^4}K_{ab}(z+nh).
\]
Its term $h=0$ is $K_{ab}(z)$. All other terms have displacement
at least $n-|z|_1$ and their sum tends to zero by absolute summability.
Thus these matrices recover every cover coefficient as $n\to\infty$.
In either case the cluster limit is unique. If the original sequence
did not converge to it in some coordinate, a subsequence staying a
fixed distance away would have a convergent further subsequence,
a contradiction. Hence the full sequence converges, with
\eqref{eq:free-horizon-narrower} at every smaller width.

\subsection{The primitive multipliers along the actual scale schedule}
Fix a scale $m$, and let $v_{d,m}$ be the train consisting of the next
$d$ prescribed blocking steps, with a top train at its fine end. Its
mass side $M_{d,m}$ tends to infinity, since each block side is at least
three. The actual path at terminal cutoff $N$ uses precisely
$v_{N-m,m}$; for fixed $m$, its depth $N-m$ tends to infinity.
Consequently each depth limit below also applies to the selected terminal
cutoff sequence. For real momentum $p$ and a direction $a$, put
\[
 u_{d,m,a}(p)=D_{v_{d,m},a}(p/M_{d,m})
                         D_{v_{d,m},a}(-p/M_{d,m}),
\]
using the cumulative drive of
Section~\ref{sec:ex-lattice-free-kernel-components}.
A single step with patch $P_m$ and side $L_m$ gives the exact identity
\[
 u_{d+1,m,a}(p)=
  |\operatorname{pullLine}_{P_m,a}(p/L_m)|^2
                          u_{d,m+1,a}(p/L_m).
\]
Each rail and centered plank is an average of unit complex numbers on
real momenta, up to a unit phase. Its modulus is at most one. Thus the
factor in front lies in $[0,1]$. Starting with $u_{0,m,a}=1$, induction
on $d$ in the displayed recurrence shows
\[
                   0\le u_{d+1,m,a}(p)\le u_{d,m,a}(p)\le1.
\]
This proves convergence for every real $p$, also after any fixed
reciprocal translation of $p$.

The noise list of $v_{d+1,m}$ is the list for $v_{d,m}$ with one
additional train appended. This follows by induction from the list
rule under one hop: both old lists undergo the same hop, with the
same top element prepended. Each additional real-momentum feed is
nonnegative. The accumulated noise $N_{d,m,a}(p)$ is therefore
nondecreasing. The uniform noise estimate of
Section~\ref{sec:ex-lattice-free-kernel-components} bounds it above,
independently of $d$ for $p$ in the base cube. At depth one it is
one. Its limit $N_{m,a}(p)$ consequently satisfies
\[
                            1\le N_{m,a}(p)<\infty.
\]

We also need the high-alias part $H_{v_{d,m},a}(p)$. Extend its
finite sum to $\mathbb Z^4$ by setting it to zero outside the centered
representatives of the nonzero residues of $M_{d,m}$. Each fixed
$z\ne0$ is represented for all sufficiently large $d$. Put
$q=p+2\pi z$. The corresponding summand is then exactly
\[
                         \frac{u_{d,m,a}(q)}{d_{M_{d,m}}(q)}.
\]
The numerator has a limit by the preceding monotonicity. The
identity
\[
 M(e^{iq_a/M}-1)\longrightarrow iq_a
\]
gives $d_M(q)\to |q|^2$. For $p$ in the base cube and $z\ne0$,
$q\ne0$, so its denominator has a nonzero limit. The uniform alias
bounds already proved give one summable majorant
\[
       C\prod_{a=1}^4(1+|z_a|)^{-2}
\]
for these extended summands. The sum of this product is finite by
the one-dimensional inverse-square series. Dominated convergence
for the series proves convergence of the high-alias part. At $z=0$
the extended summand is identically zero; no denominator at the
origin is used in this argument.

For $p\ne0$ the full directional covariance multiplier is
\[
 L_{d,m,a}(p)=
 \frac{u_{d,m,a}(p)}{d_{M_{d,m}}(p)}
       +H_{v_{d,m},a}(p)+k^{-1}N_{d,m,a}(p),\qquad k>0.
\]
It has a finite real limit $L_{m,a}(p)\ge k^{-1}>0$.
The pulse and dual-pulse symbols are the rational matrix expressions
in these multipliers given in
Section~\ref{sec:ex-lattice-free-kernel-estimates}.
Their directional denominators have just been shown to have positive
limits. The remaining scalar denominator has the form
\[
        \sum_a\frac{|e^{ip_a}-1|^2}{L_{m,a}(p)}.
\]
It is strictly positive off the origin in the base cube: every term
is nonnegative and some numerator is nonzero. Consequently every
entry of both symbols converges there. At zero their normalized
values are the identity for every cutoff. The uniform strip bounds
of the free-kernel estimates and the first subsection now give
convergence of both rows on the cover, with a retained exponential
bound.

\subsection{The adjacent one-form and two-form transports}
Fix the block side $L=L_m$. Unlike the mass side in the preceding
high-alias sum, $L$ does not change with the terminal cutoff.
For $p\ne0$ in the base cube, all the finitely many adjacent aliases
$q_l=(p+2\pi\ell_L(l))/L$ are nonzero real momenta in that cube.
Apply the directional limit above at the coarse momentum $p$ and
the fine momenta $q_l$. The principal and nonprincipal transport
formulas of Section~\ref{sec:kernel-transport-identities} then have
convergent factors.

We spell out the additional denominator check in those formulas.
The finite alias denominator is
\[
 \mathcal V_N(p)=\sum_l |m_{P_m}(q_l)|^2
                          \sum_a L_{N,a}^{\mathrm{scaled}}(q_l).
\]
Here $L_{N,a}^{\mathrm{scaled}}$ is the fine directional multiplier
with the fixed positive scaling $L^{-2}$. Every summand has a
nonnegative real limit, each inner sum has a strictly positive
limit, and the principal patch mass $m_{P_m}(p/L)$ is nonzero on
the base cube. Hence $\lim_N\mathcal V_N(p)>0$.
All other divisions in the transport formulas are by the positive
directional or scalar denominators just checked. Finite matrix
products, sums and division by these nonzero limits prove convergence
of every alias transport. Fourier reconstruction over the finite
set of $L^4$ aliases proves convergence of the adjacent transport
symbols. Their uniform row bounds again identify unique cover
limits. Possible exceptional momenta at the origin are removed by
continuity of the cluster symbols, not by taking a quotient there.

\subsection{Convergence of amendments, exterior lifts and roots}
Several operations preserve the cover convergence just established.
Sums, transposes and scalar multiples do so directly. For products,
take the minimum of two positive available widths; the convolution
of rows has a uniform bound there, multiplied by the finite middle
coordinate count. Its symbol is the product of the two symbols,
so the first subsection identifies its unique limit. For an inverse
row with a uniform bound, suppose $K_N^{-1}K_N=I$ on every torus.
Every cluster limit $Q$ of the inverse rows satisfies $QK=I$ there
by finite matrix multiplication. A square matrix with a left inverse
is invertible, so that inverse is unique. Recovery from all periods
then identifies $Q$ on the cover, and proves convergence.

The exterior lifts in Section~\ref{sec:curvature-horizon-producer}
are finite combinations of coordinate evaluation at $p_a=0$,
fixed exterior matrices and divided differences
\[
                \frac{\widehat K_N(p)-\widehat K_N(p|_{p_a=0})}
                     {e^{ip_a}-1}.
\]
Their uniform bounds at a smaller positive width were proved there.
For $p_a\ne0$ the displayed expression converges, because both
numerator terms do and the denominator is fixed. Thus any two
cluster symbols agree off that coordinate hyperplane, and by
continuity everywhere. This proves convergence of the divided
row. Repeating the argument through the finitely many axes proves
convergence of the one-sided and two-sided exterior lifts, with
all their fixed matrix factors.

It remains to justify the positive square roots; convergence of an
arbitrary matrix factor would be false. Suppose $P_N$ have a common
exponential row bound, their torus matrices are positive semidefinite,
and $P_N^{\mathsf T}P_N$ converge on every torus. Two cluster limits
$P,Q$ have positive semidefinite torus matrices, since symmetry and
nonnegativity of each quadratic form pass to the limit. Their
squares are equal. A positive semidefinite square root is unique:
if $A\ge0$ and $A^2=B$, then $A$ commutes with $B$; on each
eigenspace of $B$, diagonalizing the symmetric restriction of $A$
forces all eigenvalues to be the nonnegative square root of that
eigenvalue of $B$. Hence $A$ is fixed on every eigenspace.
Apply this to $P$ and $Q$ on each torus, then recover the cover
coefficients. The unique-cluster-limit argument proves convergence
of $P_N$ itself.

The amended beds have exactly these properties. If $B_N$ is the
dual-pulse row, its amended inverse precision is
\[
 A_N=B_N+l(d_2B_N)^{\mathsf T}(d_2B_N),\qquad l\ge0,
\]
and its precision is $G_N=A_N^{-1}$. Products and the inverse argument
give their limits. The root row has square $G_N-cI$, where $c$ is
the fixed direct clipping coefficient. Its uniform strip bound
and positive semidefinite torus values are supplied by
Section~\ref{sec:positive-row-producer}, so the positive-root argument
applies. The original pulse, dual pulse, amended precision, its
inverse and the root row therefore converge together.

For the curvature root, use the explicit amended row of
Section~\ref{sec:curvature-horizon-producer}:
\[
 R_N=D_N^{\mathsf T}(C_N+lI-hS_N^{\mathsf T}S_N)D_N,
 \qquad D_N=(I+l d_2 B_{c,N}d_2^{\mathsf T})^{-1}.
\]
Here $l,h$ are the fixed coarse and fine amendment coefficients;
$C_N,S_N$ are the two exterior lifts constructed there. Their
convergence follows from the operations just proved. The positivity
and uniform strip price of $R_N$ and its positive root are already
established. Since that root's square is $R_N$, it too converges.
All of these statements concern the actual chosen positive roots;
no new choice of signs or bases has been made.

\subsection{Completing the retained tower}
The full precision factors $J_{f,N},J_{c,N}$ consist of the fixed
direct block and the root rows, so they converge on every torus.
Write $G_{f,N},G_{c,N}$ for the corresponding precisions and
$R_{f,N}=G_{f,N}^{-1}$, $R_{c,N}=G_{c,N}^{-1}$. With the two-form
transport $T_N$ and curvature root $P_N$, the exact tower blocks are
\begin{align*}
 M_{F,N}&=J_{f,N}R_{f,N}T_NJ_{c,N}^{\mathsf T},\\
 M_{N,N}&=k^{-1/2}d_{1,c}^{\mathsf T}J_{c,N}^{\mathsf T},\\
 M_{O,N}&=
 \begin{bmatrix}
 I-J_{c,N}R_{c,N}J_{c,N}^{\mathsf T}\\
 P_Nd_{2,c}R_{c,N}J_{c,N}^{\mathsf T}
 \end{bmatrix}.
\end{align*}
These are the same formulas as
\eqref{eq:ex-lattice-horizon-blocks}, on each selected physical torus
of side $n_m(p)$, with its prescribed fine partner. Here $n_m(p)$
denotes the period sequence fixed at scale $m$ by the schedule.
Here $p$ is fixed while $N$ increases. The committed transition
identities apply on precisely these tori. Every finite matrix factor
has a limit there, hence so does each block. Their
uniform strip bounds were purchased in the retained tower
construction. The period-recovery criterion from the first
subsection gives convergence of their cover coefficients. Use the selected sides above, which tend to infinity as $p$ does;
no assumption on convergence of an unrelated periodic realization
is needed.

Together with the one-form predictor and two-form transport, these
are precisely the five retained tower rows. Thus at every fixed
scale all five bed rows and all five tower rows converge, and each
family keeps some positive exponential width and finite budget.
The purchases of widths and budgets are made at the fixed stages,
before choosing the later terminal cutoff. The eventual
normalization changes discussed in
Section~\ref{sec:selected-history-limits} preserve these limits.
This supplies the free-horizon convergence input of that section
from the actual recursively selected free kernels.
